\clearpage
\section{Attempts to widen the zones}
\label{app:widen}

Five attempts to make the \cert{} zones wider without more wrong certificates (the zones certify \WidenRefTfivePLow{}\% to \WidenRefTfivePHigh{}\% of machine-paraphrased inputs for their own task, by pivot; Section~\ref{sec:paraphrase}). Each had its setting and success line written in the plan of the repository before the run. All work on routes only, on a CPU, and train nothing. One split throughout: \WidenTasksCal{} calibration tasks, in halves A and B of \WidenHalfA{} and \WidenHalfB{} where a study needs two, and \WidenTasksEval{} evaluation tasks in one random order of arrival. \textsf{C}, \textsf{U}, \textsf{P}, \textsf{M}, \textsf{R} are defined in Section~\ref{sec:paraphrase}; \textsf{Z} is the share of paraphrased inputs inside their own task's zone, whatever an older zone does with them. Theorem~\ref{thm:writeonce} holds for every zone below as defined, since each is a fixed set once its task has arrived, decided by frozen parts computed reproducibly (Remark~\ref{rem:otherzones}); for the judge zone as measured it does not (idea 4). The machine paraphrases are not rewordings by people, one radius for all tasks is used throughout, and a much larger judge (for example the frozen base model asked whether two instructions ask for the same thing) was not tried. We do not claim that no zone can do better, only that none of these five met its line.

\paragraph{Idea 1: hidden states of a frozen causal language model.} \emph{Tried.} For a decoder-only model the hidden states at the first $n$ positions depend only on the first $n$ tokens, so one pass gives a feature in the model's own reading; zones, radii, routes and the theorem are unchanged. SmolLM2-135M, layers 0, 2, 4, 8, 16, 30, pooled as a weighted mean over the first $n$ positions or as the state at position $n$ (\WidenFeatSettings{} settings). Layer, pooling and radius (\cert{} with $\kappa = 0.1, \dots, 1.0$, or one radius at a quantile of 0.5 to 0.95) were chosen on the calibration tasks for the highest \textsf{P} among settings with \textsf{C} at most that of \cert{} there (\WidenFeatCalRefC{}\%), then used once. Layer 0 with the weighted mean reproduces the routes published for this model, as it had to.
\emph{Line fixed before the run:} success if at least 20 points more inputs are certified for their own task than with word-vectors at the same or a lower collision rate and no more wrong certificates; against if no setting gains 5 points at a matched collision rate or wrong certificates go up. \emph{Verdict: against} (Table~\ref{tab:widenmain}).
\emph{Result.} Chosen: layer 30, weighted mean, one radius at quantile \WidenFeatRadiusQuantile{}. On the evaluation tasks \textsf{P} rises to \WidenFeatP{}\% (\WidenFeatGain{} points over the reference row), but \textsf{C} and \textsf{M} rise too (Table~\ref{tab:widenmain}). The collision rate of the calibration tasks (\WidenFeatCalC{}\% for the chosen setting) does not carry over to \WidenTasksEval{} tasks, since a longer sequence has more collisions, so the first clause could not be applied as written; the verdict is not softened for that.
\emph{Learned.} Hidden states tolerate paraphrase better and tell siblings apart worse; they move along the trade and do not leave it.

\paragraph{Idea 2: a metric in units of what paraphrasing does.} \emph{Tried.} The word-vector feature and zones stay; the ruler changes. Instead of rescaling by how unrelated inputs vary, distances are rescaled by the second moment of ``feature of the paraphrase minus feature of the original'', shrunk by 10\% towards the identity; the feature map is fixed after set-up, so the theorem holds. Moments from the first three inputs of the \WidenTasksCal{} calibration tasks, German pivot only; T5-base and SmolLM2 word-vectors; evaluation on the French pivot, which the ruler has not seen; one radius swept over a grid; old and new ruler read at the radius where \textsf{C} equals that of \cert{} (the \textsf{C} of \cert{} on T5-base).
\emph{Line fixed before the run:} success if at that point \textsf{P} gains at least 20 points with no more \textsf{M}; against if the gain is below 5 points or \textsf{M} rises. \emph{Verdict: against} for both models.
\emph{Result.} T5-base: gain in \textsf{P} \WidenRulerTfiveGain{} points, \textsf{M} rises (Table~\ref{tab:widenmain}). SmolLM2: gain \WidenRulerSmolGain{}, \textsf{M} from \WidenRulerSmolMOld{}\% to \WidenRulerSmolMNew{}\%. With the paper's own radius rule the gain on T5-base is nil (\textsf{P} \WidenRulerTfiveRuleBOld{}\% against \WidenRulerTfiveRuleBNew{}\%, French pivot); with 50\% shrinkage \textsf{P} is \WidenRulerHalfP{}\% and \textsf{M} \WidenRulerHalfM{}\%, which changes no verdict. On the five benchmark tasks with the development rewordings the new ruler is worse: with the heavy rewording \WidenRulerHeavyNew{}\% of inputs reach their own task, against \WidenRulerHeavyOld{}\%.
\emph{Learned.} The ten largest directions of the movement hold only \WidenRulerTopTenTfive{}\% (T5-base) and \WidenRulerTopTenSmol{}\% (SmolLM2) of it, so it is not confined to a few directions, which may be why a linear ruler on an average of word-vectors did not separate it from the difference between siblings.

\paragraph{Idea 3: a frozen sentence encoder and a radius with a stated coverage.} \emph{Tried.} A sentence encoder, made to bring paraphrases together, as the feature, with one radius set by split-conformal calibration so that a share $1 - \alpha$ of rewordings lies inside ($\alpha$ is the level of the radius, not the scaling factor of the adapters; Proposition~\ref{prop:conformal}). all-MiniLM-L6-v2 (\WidenEncParams{} million weights), mean pooling, unit length, Euclidean distance, on the first tokens of the input. The opening length $n_k$ of a task, in encoder tokens, is rounded down to the grid \WidenEncGrid{} (\WidenEncPrefixes{} of these lengths occur), so an unedited input is at distance exactly 0 from its own signature (true for all \WidenEncUneditedZero{} unedited inputs). The ruler, if used, is fitted on half A (\WidenEncDisplacements{} displacements, German pivot, 10\% shrinkage); the radius is the $\lceil (1-\alpha)(m+1) \rceil$-th smallest distance of the first paraphrased input of each task of half B, German pivot ($m = \WidenHalfB{}$), for $\alpha = 0.1$ (primary), $0.05$, $0.2$. Plain or ruler as primary was decided on half B only, by the larger ratio of the median distance between a signature and the nearest other signature to the radius: \WidenEncRatioRuler{} for the ruler against \WidenEncRatioPlain{}, so the ruler. Routes as in \cert{}.
\emph{Line fixed before the run:} success if at $\alpha = 0.1$ on the French pivot \textsf{P} $\ge$ 75\%, \textsf{M} $\le$ 5\% and \textsf{C} $\le$ 19\%; against if \textsf{P} $<$ 55\%, \textsf{M} $>$ 9\% or \textsf{C} $>$ 25\%; else partial. \emph{Verdict: against} (\textsf{M} above 9\%, \textsf{P} below 75\%).
\emph{Result} (Table~\ref{tab:widenencoder}). Primary setting: \textsf{P}~=~\WidenEncPrimP{}\%, \textsf{M}~=~\WidenEncPrimM{}\%. Over the six settings \textsf{P} rises from \WidenRefTfiveFrP{}\% (\cert{}, French) to between \WidenEncRulerTwentyFrP{}\% and \WidenEncPlainTenFrP{}\% and \textsf{R} rises as well, but \textsf{M} is at least \WidenEncRulerTwentyFrM{}\% in all six, and \WidenEncPlainTenFrM{}\% where \textsf{P} is highest. A second encoder (bge-small-en-v1.5) gives the same picture (\textsf{P}~\WidenBgeP{}\%, \textsf{M}~\WidenBgeM{}\%, \textsf{C}~\WidenBgeC{}\%). With the instruction as a separate field, so that the whole definition is encoded, three encoders (all-MiniLM-L6-v2, bge-small-en-v1.5 \citep{xiao2023cpack}, all-mpnet-base-v2 \citep{song2020mpnet}) give \textsf{P} from \WidenFieldPLow{}\% to \WidenFieldPHigh{}\% and \textsf{M} from \WidenFieldMLow{}\% to \WidenFieldMHigh{}\%. The second planned view, the curve over a grid of radii read at the collision rate of \cert{}, was run at \textsf{C}~=~\WidenEncCurveC{}\% (the SmolLM2 reference, not T5-base): ruler \textsf{P}~\WidenEncCurveRulerP{}\%, \textsf{M}~\WidenEncCurveRulerM{}\%, \textsf{R}~\WidenEncCurveRulerR{}\%; plain \WidenEncCurvePlainP{}\%, \WidenEncCurvePlainM{}\%, \WidenEncCurvePlainR{}\%. The ruler at $\alpha = 0.2$ was looked at after the run, not chosen, and does not meet the line; on the French pivot it has fewer wrong and more correct certificates than \cert{} (\textsf{M} \WidenRefTfiveFrM{}\%, \textsf{P} \WidenRefTfiveFrP{}\% for \cert{}), on the German pivot more wrong ones (\WidenEncRulerTwentyDeM{}\% against \WidenRefTfiveDeM{}\%).
\emph{Learned.} The encoder does what it was chosen for but does not bring wrong certificates under the bar.

\begin{table}[t]
  \caption{The six settings of the sentence encoder (plain metric or paraphrase ruler; $\alpha = 0.05, 0.1, 0.2$) on the French and German pivot, \WidenParaInputsPivot{} paraphrased inputs each, \WidenTasksEval{} evaluation tasks, in \%. Primary setting in bold. \textsf{C} and \textsf{U} do not depend on the pivot. \textsf{Z} is to be compared with $1 - \alpha$. The radius was set on German-pivot paraphrases of other tasks.}
  \label{tab:widenencoder}
  \centering
  \footnotesize
  \setlength{\tabcolsep}{4pt}
  \begin{tabular}{llrrrrrrr}
    \toprule
    Metric & $\alpha$ & Pivot & \shortstack[r]{tasks that\\collide\\(\textsf{C}), \% $\downarrow$} & \shortstack[r]{own task,\\unedited\\(\textsf{U}), \% $\uparrow$} & \shortstack[r]{right\\certificate\\(\textsf{P}), \% $\uparrow$} & \shortstack[r]{wrong\\certificate\\(\textsf{M}), \% $\downarrow$} & \shortstack[r]{own task,\\paraphrased\\(\textsf{R}), \% $\uparrow$} & \shortstack[r]{inside own\\zone\\(\textsf{Z}), \% $\uparrow$} \\
    \midrule
    plain & 0.05 & French & 39.0 & 62.0 & 65.5 & 29.8 & 69.7 & 95.1 \\
     &  & German & 39.0 & 62.0 & 67.8 & 29.7 & 69.3 & 97.3 \\
    \addlinespace[2pt]
     & 0.1 & French & 22.9 & 77.6 & 73.1 & 17.1 & 82.0 & 89.8 \\
     &  & German & 22.9 & 77.6 & 75.0 & 17.2 & 81.5 & 92.0 \\
    \addlinespace[2pt]
     & 0.2 & French & 15.7 & 84.4 & 69.1 & 10.1 & 87.5 & 78.9 \\
     &  & German & 15.7 & 84.4 & 69.0 & 10.1 & 88.1 & 79.1 \\
    \addlinespace[2pt]
    \midrule
    ruler & 0.05 & French & 29.9 & 70.7 & 69.1 & 24.1 & 75.0 & 92.7 \\
     &  & German & 29.9 & 70.7 & 67.4 & 29.1 & 69.9 & 96.3 \\
    \addlinespace[2pt]
     & 0.1 & French & \textbf{16.6} & \textbf{83.5} & \textbf{71.2} & \textbf{11.6} & \textbf{86.1} & \textbf{82.0} \\
     &  & German & 16.6 & 83.5 & 73.9 & 14.5 & 83.5 & 88.3 \\
    \addlinespace[2pt]
     & 0.2 & French & 12.1 & 88.1 & 62.1 & 6.0 & 90.0 & 68.0 \\
     &  & German & 12.1 & 88.1 & 72.6 & 8.1 & 89.7 & 80.4 \\
    \bottomrule
\end{tabular}

\end{table}

\paragraph{Idea 4: a zone decided by a frozen judge of ``the same instruction''.} \emph{Tried.} In every measure of likeness a sibling task is as close as a paraphrase, while the question a zone must answer is whether two instructions ask for the same thing. The theorem needs no distance, only a set fixed when the task arrives. The instruction is a separate field. The zone of task $k$ holds the instructions $u$ with (a) MiniLM distance to the task's instruction at most $r_w$ and (b) judge score at least $\tau$: the smaller of the two probabilities of ``entailment'' ($u$ given the instruction and the reverse) from cross-encoder/nli-deberta-v3-xsmall \citep{he2023debertav3}, texts cut at 128 tokens, 1 for identical texts. Thresholds from half B, German pivot, by split-conformal quantiles ($r_w$ at $\alpha = 0.02$, $\tau$ at $\alpha = 0.08$, so that at least 90\% of rewordings are inside). A free judge (largest word weight among words in only one text) was to be used if better by more than 0.02 in area under the curve; it was not (\WidenJudgeAucFree{} against \WidenJudgeAucNli{}). Route: the oldest task whose zone holds the instruction, else the nearest by the encoder. At most \WidenJudgeLimit{} nearest tasks inside $r_w$ are judged per text. This limit is a shortcut of the evaluation: the \WidenJudgeLimit{} are the nearest of all \WidenTasksEval{} tasks, older and newer, so the zone as measured also depends on later tasks and Theorem~\ref{thm:writeonce} does not apply to it; it applies to the zone as defined, without a limit, which can only claim more texts; we did not measure that.
\emph{Line fixed before the run:} the bar of idea 3. \emph{Verdict: against} (\textsf{M} far above 9\%, \textsf{C} above 25\%, \textsf{P} below 75\%; Table~\ref{tab:widenmain}).
\emph{Result.} $r_w$ came out as \WidenJudgeRw{}, the largest calibration distance (with \WidenHalfB{} tasks the $\alpha = 0.02$ quantile is the maximum), so the filter lets almost every pair through and alone would certify almost every text for a wrong task (\textsf{P} \WidenJudgeFilterOnlyP{}\%, \textsf{M} \WidenJudgeFilterOnlyM{}\%). With $\tau = \WidenJudgeTau{}$, \WidenJudgePosSame{}\% of calibration paraphrase pairs and \WidenJudgeNegSame{}\% of pairs of nearest other tasks count as the same request. Evaluation, French pivot (Table~\ref{tab:widenmain}): \textsf{M}~=~\WidenJudgeFrM{}\%, \textsf{Z}~=~\WidenJudgeFrZ{}\%; German: \textsf{P}~\WidenJudgeDeP{}\%, \textsf{M}~\WidenJudgeDeM{}\%, \textsf{R}~\WidenJudgeDeR{}\%. A stricter threshold does not rescue it (Table~\ref{tab:widenjudge}, read after the run, nothing chosen): \textsf{P} never exceeds \WidenSweepPMax{}\%; \textsf{M} is below 9\% only at the strictest threshold, where almost nothing is certified; where \textsf{C} is nearest to 19\% ($\tau = \WidenSweepMidTau{}$) the result is close to the encoder alone. Of the \WidenJudgeCollisionPairs{} pairs of tasks that collide under \cert{}, the judge calls \WidenJudgeCollisionSame{} the same (\WidenJudgeCollisionNotSameSame{} of the \WidenJudgeCollisionNotSame{} whose definitions are not identical). On the five benchmark tasks it handles free rewording better than any distance: \WidenJudgeBenchInside{} of the \WidenJudgeBenchTotal{} development rewordings lie inside their own zone (not the heavy rewording of the sentiment task), all go to their own task, none is certified for another.
\emph{Learned.} The judge is good on single pairs (area under the curve \WidenJudgeAucNli{} on half B) and bad as a zone: the limit was reached for \WidenJudgeLimitFr{} of the \WidenTasksEval{} French texts, and one wrong ``same'' among up to \WidenJudgeLimit{} nearby tasks, from an older zone or when the text's own zone does not hold it, gives a wrong certificate: \WidenJudgeNegSame{}\% wrong ``same'' on the nearest pair becomes \WidenJudgeFrM{}\% wrong certificates.

\begin{table}[t]
  \caption{The judge zone over twelve thresholds $\tau$ (French pivot, \WidenTasksEval{} evaluation tasks, in \%), looked at after the run; nothing was chosen from it. The primary $\tau = \WidenJudgeTau{}$ was fixed on half B before the run (the row 0.19 is a grid point).}
  \label{tab:widenjudge}
  \centering
  \small
  \begin{tabular}{rrrrrr}
    \toprule
    $\tau$ & \shortstack[r]{tasks that\\collide\\(\textsf{C}), \% $\downarrow$} & \shortstack[r]{inside own\\zone\\(\textsf{Z}), \% $\uparrow$} & \shortstack[r]{right\\certificate\\(\textsf{P}), \% $\uparrow$} & \shortstack[r]{wrong\\certificate\\(\textsf{M}), \% $\downarrow$} & \shortstack[r]{own task,\\paraphrased\\(\textsf{R}), \% $\uparrow$} \\
    \midrule
    0.01 & 78.9 & 96.6 & 22.3 & 77.1 & 22.9 \\
    0.10 & 46.4 & 92.9 & 49.4 & 46.3 & 53.6 \\
    0.19 & 36.3 & 90.6 & 58.6 & 35.0 & 64.4 \\
    0.28 & 31.1 & 89.0 & 62.6 & 29.7 & 69.7 \\
    0.37 & 27.3 & 87.9 & 64.7 & 26.1 & 73.3 \\
    0.46 & 25.6 & 87.4 & 67.0 & 23.6 & 75.9 \\
    0.55 & 23.4 & 86.6 & 68.6 & 21.1 & 78.3 \\
    0.64 & 20.9 & 85.6 & 68.7 & 20.1 & 79.3 \\
    0.73 & 19.0 & 84.0 & 69.3 & 17.7 & 81.3 \\
    0.82 & 14.4 & 82.1 & 69.9 & 14.7 & 84.1 \\
    0.91 & 10.4 & 79.7 & 70.4 & 10.6 & 88.3 \\
    1.00 & 1.4 & 0.6 & 0.6 & 0.0 & 96.4 \\
    \bottomrule
\end{tabular}

\end{table}

\paragraph{Exploration, not a test: why likeness is not enough.} On the \WidenTasksCal{} calibration tasks and their German-pivot paraphrases only (nothing from the evaluation tasks; numbers copied from the notes of the repository, the script was not kept): a paraphrase drops \WidenExploreDropShare{}\% of the words of a definition, and rare words more often than common ones (\WidenExploreDropRare{}\% of words in only one calibration definition, \WidenExploreDropCommon{}\% of words in more than 60). Sibling tasks (\WidenExploreSiblingPairs{} pairs with different texts sharing at least \WidenExploreSiblingSharedMin{}\% of their words) differ in the same kind of words: \WidenExploreSiblingRare{}\% of the words in which they differ occur in only one definition, \WidenExploreSiblingFew{}\% in at most three. So ``the rare words must all be present'' fails: only \WidenExploreKeepAllOne{}\% of paraphrases keep all words that occur in one definition only (\WidenExploreKeepAllTwo{}\% for at most two), and \WidenExploreSiblingSameOne{}\% (\WidenExploreSiblingSameTwo{}\%) of sibling pairs have the same such words. Rewording and being a sibling change the same words, so no word count, rarity weighting or distance between averages can tell them apart; ideas built on key words or on learned weights were dropped without a test.
\emph{A second exploration, a round trip as a normal form} (calibration tasks): passing the original and its paraphrase through a frozen translator (English to German and back) made \WidenExploreRoundTripIdentical{}\% of the sentences of a French-pivot paraphrase identical to those of the original's own round trip, and raised the share of sibling comparisons in which the sibling is farther from the original than the own paraphrase is, in a word-level measure, from \WidenExploreRoundTripFarFrom{}\% to \WidenExploreRoundTripFarTo{}\%. It was not tested as a zone.

\paragraph{Idea 5: a zone decided sentence by sentence.} \emph{Tried.} Idea 4 asked the judge about whole instructions of up to 128 tokens, where one decisive word is easily lost, and against up to \WidenJudgeLimit{} candidates at once; here the judge is asked sentence by sentence and only about close candidates. The instruction is a separate field. The zone of task $k$ holds the texts $u$ with (1) MiniLM distance (whole text) to the stored instruction $o_k$ at most $r$ and (2) sentence-by-sentence score at least $\tau$: $u$ and $o_k$ are split into sentences, each sentence of $u$ is paired with its most similar sentence of $o_k$ and the reverse (MiniLM cosine); identical sentences score 1, any other pair the smaller of the two ``entailment'' probabilities of a frozen inference model; the score of $(u, o_k)$ is the smallest over all pairs. Every task within $r$ is judged, with no limit, so the zone is a fixed set and Theorem~\ref{thm:writeonce} applies to the zone as measured. A text outside every zone goes to the task with the nearest MiniLM signature (no certificate). Thresholds from half B, German pivot, by split-conformal quantiles: $r$ at $\alpha = \WidenSentRAlpha{}$ ($r = \WidenSentR{}$), $\tau$ at $\alpha = \WidenSentTauAlpha{}$. The judge was chosen on half B only, among three, by the area under the curve for telling a task's German paraphrase (\WidenSentPositives{} pairs) from the other calibration instructions within $r$ (\WidenSentNegatives{} pairs): \WidenSentAucXsmall{} for cross-encoder/nli-deberta-v3-xsmall (used, also the smallest), \WidenSentAucBase{} for nli-deberta-v3-base and \WidenSentAucMnli{} for MoritzLaurer/DeBERTa-v3-base-mnli-fever-anli.
\emph{Line fixed before the run:} the bar of ideas 3 and 4 (success: \textsf{P} $\ge$ 75\%, \textsf{M} $\le$ 5\%, \textsf{C} $\le$ 19\%; against: \textsf{P} $<$ 55\%, \textsf{M} $>$ 9\% or \textsf{C} $>$ 25\%), read on the Italian pivot, which no study had read. \emph{Verdict: against}: on the Italian pivot \textsf{M} is \WidenSentItM{}\%, above 9\%, while \textsf{P} (\WidenSentItP{}\%) and \textsf{C} (\WidenSentC{}\%) are on the right side of the success lines. Expected: \textsf{P} 62\% to 74\%, \textsf{M} 3\% to 7\%, \textsf{C} 5\% to 10\%; \textsf{P} came out above, \textsf{M} and \textsf{C} above theirs.
\emph{Calibration.} The threshold is very low, $\tau = \WidenSentTau{}$: the score is the smallest over all sentence pairs, so a true rewording of one sentence often scores low on some pair, and keeping \WidenSentPosSameXsmall{}\% of German paraphrases takes the threshold close to zero. There \WidenSentNegSameXsmall{}\% of pairs of nearest other instructions still count as the same (\WidenSentNegSameBase{}\% and \WidenSentNegSameMnli{}\% for the other two judges at their own thresholds): the test hardly filters beyond what the encoder filter lets through.
\emph{Result} (Table~\ref{tab:widenmain}). \textsf{C} is the same on every pivot. Italian \textsf{Z} \WidenSentItZ{}\%; French \WidenSentFrP{}\%, \WidenSentFrM{}\%, \WidenSentFrR{}\%, \WidenSentFrZ{}\%; German \WidenSentDeP{}\%, \WidenSentDeM{}\%, \WidenSentDeR{}\%, \WidenSentDeZ{}\%; Spanish \WidenSentEsP{}\%, \WidenSentEsM{}\%, \WidenSentEsR{}\%, \WidenSentEsZ{}\%. It is the best point: fewer wrong certificates and about half the collisions of the sentence encoder alone (Table~\ref{tab:widenencoder}). The wrong certificates are the usual siblings, for example \texttt{\WidenSentExOneTask{}} claimed by \texttt{\WidenSentExOneClaim{}} (lowest sentence-pair score \WidenSentExOneScore{}), \texttt{\WidenSentExTwoTask{}} by \texttt{\WidenSentExTwoClaim{}} (\WidenSentExTwoScore{}) and \texttt{\WidenSentExThreeTask{}} by \texttt{\WidenSentExThreeClaim{}} (\WidenSentExThreeScore{}), the first three of the \WidenSentWrongCerts{} wrong certificates on the Italian pivot in the file. Of the \WidenSentCollNotSame{} colliding pairs of tasks with different texts the judge still passes \WidenSentCollPass{}. On the five benchmark tasks the development rewordings are mostly outside these zones (\WidenSentBenchLightZone{} of \WidenSentBenchN{} light, \WidenSentBenchMediumZone{} medium, \WidenSentBenchHeavyZone{} heavy lie inside), all are routed to the right task, none is certified for a wrong one. Cost: \WidenSentSeconds{} s per text on \WidenSentThreads{} CPU threads; a text has \WidenSentWithinMean{} tasks within $r$ on average (\WidenSentMax{} at most over the four pivots); \WidenSentNoneIt{} of the \WidenTasksEval{} Italian texts have none and are not certified.
\emph{Learned.} The frontier moved (about 40\% certified with about 9\% wrong for \cert{}, \WidenSentItP{}\% with \WidenSentItM{}\% here) and the wall did not: what separates a rewording from a sibling is the meaning of one or two words, and no small frozen model we tried knows it (a small inference model is not shown to know that ``reaction'' and ``intent'' ask for different things; we did not test why).

\paragraph{Sibling tasks that collide.} Table~\ref{tab:widensiblings} lists pairs from the files of the runs, selected by hand. They differ in one or a few decisive words (a date format, an order of comparison, which part of a relation is asked for), and a feature that reads meaning, or a judge, sees nearly the same request. The files also hold wrong certificates between tasks that are not siblings: sentence generation from a relation for an entailment task (judge score 0.47), option generation for a physics-style question for a gap classification task (0.68), gender criticism for offensiveness (0.74).

\begin{table}[t]
  \caption{Sibling tasks that collide. (3): the unedited instruction of the arriving task lies inside the zone of the older task under the encoder zone (first 15 collisions of the file). (4): a French paraphrase of the first task is certified for the second (older or, when its own zone does not hold it, newer) under the judge zone (first 12 such cases); the score is the judge's score for that paraphrase against the instruction of the claiming task. ``--'': not among the cases listed.}
  \label{tab:widensiblings}
  \centering
  \footnotesize
  \setlength{\tabcolsep}{3pt}
  \scriptsize\begin{tabular}{llcc}
    \toprule
    Task that arrives & Task whose zone claims it & Listed under & Judge score \\
    \midrule
    \texttt{check\_validity\_date\_mmddyyyy} & \texttt{check\_validity\_date\_ddmmyyyy} & (3), (4) & 0.46 \\
    \texttt{farthest\_integers} & \texttt{closest\_integers} & (3) & -- \\
    \texttt{atomic\_classification\_xintent} & \texttt{atomic\_classification\_xreact} & (3), (4) & 0.91 \\
    \texttt{all\_elements\_except\_first\_i} & \texttt{first\_i\_elements} & (3), (4) & 0.89 \\
    \texttt{synthetic\_remove\_floats} & \texttt{synthetic\_remove\_odds} & (4) & 0.65 \\
    \texttt{amazonreview\_category\_classification} & \texttt{amazonreview\_polarity\_classification} & (4) & 0.49 \\
    \bottomrule
\end{tabular}

\end{table}

\begin{proposition}[Coverage of the split-conformal radius]
\label{prop:conformal}
Let $D_1, \dots, D_m$ be the distances of the calibration rewordings from their own signatures and
$D_{m+1}$ that of a new task's rewording. If $(D_1, \dots, D_{m+1})$ are exchangeable, and $r_\alpha$ is the
$\lceil (1-\alpha)(m+1) \rceil$-th smallest of $D_1, \dots, D_m$ ($r_\alpha = \infty$ if that index exceeds $m$), then
$\Pr[D_{m+1} \le r_\alpha] \ge 1 - \alpha$, where the probability is over the joint draw of all $m+1$ values; it is not conditional on the calibration set.
\end{proposition}

\begin{proof}
Put $k=\lceil(1-\alpha)(m+1)\rceil$; if $k>m$ there is nothing to prove. Let $N_i$ be the number of $j\le m+1$ with $D_j<D_i$. If $D_{m+1}>r_\alpha$, then
$D_{m+1}$ exceeds the $k$ smallest of $D_1,\dots,D_m$, so $N_{m+1}\ge k$. Any index with $N_i\ge k$ has $D_i$ larger than the $k$-th smallest of all $m+1$ values,
so at most $m+1-k$ indices have $N_i\ge k$. By exchangeability $\Pr[N_{m+1}\ge k]=\frac{1}{m+1}\mathbb{E}\,|\{i: N_i\ge k\}|\le\frac{m+1-k}{m+1}\le\alpha$.
\end{proof}

This is the standard result of split-conformal prediction \citep{vovk2005algorithmic,angelopoulos2023conformal}. It says nothing about the zones of older tasks, which may claim the input first, nor about the route. \emph{Check} (\textsf{Z} in Table~\ref{tab:widenencoder}). With the plain metric \textsf{Z} is slightly below the promise in \WidenPlainZBelowCount{} of six cases (by at most \WidenPlainZBelowMax{} points) and above it in the others. The condition does not hold exactly (calibration paraphrases go through German, test paraphrases of the primary column through French, and the tasks differ), so the agreement is an observation. With the ruler, fitted on German-pivot paraphrases, the promise did not hold on the French pivot at any $\alpha$ (95\%, 90\% and 80\% promised), and it also misses on the German pivot at $\alpha = 0.1$ (90\% promised). For the judge zone, \textsf{Z} is \WidenJudgeFrZ{}\% on the French pivot where 90\% was promised.

\paragraph{Cost per input.} The word-vector zones cost one average and the distances to the signatures. The sentence encoder (\WidenEncParams{} million weights) runs on up to \WidenEncPrefixes{} prefixes per input: \WidenEncMs{} ms on average (\WidenEncMsHigh{} ms at the 90th percentile) on two CPU threads, including the distances to \WidenTasksEval{} signatures. The judge zone adds an inference model: about \WidenJudgePairs{} judged pairs and \WidenJudgeSeconds{} seconds per text on two threads, growing with the number of look-alike tasks. Both give up the edit bound of Proposition~\ref{prop:edits}.

\paragraph{A check on held-out paraphrases (not one of the five ideas).} \emph{Tried.} In the table of idea 3 the paraphrase ruler at $\alpha = 0.2$, not the primary setting, was better than \cert{} on all four measures on the French pivot (\textsf{C} \WidenSpC{}\% against \WidenSpCRule{}\%). A comparison chosen after the results proves nothing, so the paraphrases through Spanish (MarianMT, as above), read by no study, were used on the same \WidenTasksEval{} tasks in the same order, with the same calibration (ruler from half A, radius from half B at $\alpha = 0.2$, both from German-pivot paraphrases) and nothing fitted or chosen again; the script first reproduced the stored radius, the stored French and German numbers of the encoder setting and the numbers of \cert{}.
\emph{Line fixed before the Spanish paraphrases were read:} confirmed if the encoder setting has \textsf{P} at least \WidenSpAskP{} points and \textsf{R} at least \WidenSpAskR{} points above \cert{}, and \textsf{M} not above. \emph{Verdict: not confirmed.} We had expected \textsf{P} about 40\% for \cert{}; it is \WidenRefTfiveEsP{}\%, and these paraphrases keep more of the original words (\WidenSpKept{}\% against \WidenKeptDe{}\% and \WidenKeptFr{}\%; \WidenSpSame{} of the \WidenTasksAll{} are identical to the original, against \WidenSameFr{} French and \WidenSameDe{} German), which may explain it.
\emph{Result} (Table~\ref{tab:widenspanish}). Spanish pivot: the encoder setting has \textsf{Z} \WidenSpEncEsZ{}\%. The gain in \textsf{P} is \WidenSpGainP{} points where \WidenSpAskP{} were asked; the gain in \textsf{R} (\WidenSpGainR{}) and the condition on \textsf{M} hold. The promise of $1 - \alpha = \WidenSpPromise{}\%$ does not carry over for the ruler: \textsf{Z} is below it on Spanish, French and the German pivot it was fitted on (Table~\ref{tab:widenencoder}). The other five encoder settings repeat the French picture on the Spanish pivot (ruler at $\alpha = 0.1$: \textsf{P} \WidenSpAllRulerTenP{}\%, \textsf{M} \WidenSpAllRulerTenM{}\%; plain at $\alpha = 0.2$: \WidenSpAllPlainTwentyP{}\% and \WidenSpAllPlainTwentyM{}\%). On the five benchmark tasks the original and light rewording lie inside their own zones; the medium and heavy do not (\WidenSpBenchMediumZone{}\% and \WidenSpBenchHeavyZone{}\% inside), and all are still given to the right task (\WidenSpBenchMediumRouted{}\% and \WidenSpBenchHeavyRouted{}\%), none certified for a wrong one.
\emph{Other splits} (Table~\ref{tab:widensplits}; \WidenSplitCount{} random splits, seeds 1 to 5, ruler and radius fitted again on each; Spanish pivot). Mean \textsf{P} is \WidenSplitRulePMean{}\% for \cert{} and \WidenSplitEncPMean{}\% for the encoder setting, mean \textsf{M} \WidenSplitRuleMMean{}\% and \WidenSplitEncMMean{}\%; the encoder setting varies more between splits. No split passes the rule by itself. On \WidenSplitMoreM{} splits it has more wrong certificates than \cert{} (\WidenSplitSeedOneEncM{}\% against \WidenSplitSeedOneRuleM{}\% on seed 1, \WidenSplitSeedThreeEncM{}\% against \WidenSplitSeedThreeRuleM{}\% on seed 3, \WidenSplitSeedFourEncM{}\% against \WidenSplitSeedFourRuleM{}\% on seed 4) and on \WidenSplitMoreC{} more collisions (seed 3: \WidenSplitSeedThreeEncC{}\% against \WidenSplitSeedThreeRuleC{}\%). French pivot, same splits: \cert{} \textsf{P} \WidenSplitFrRulePMean{}\% (\WidenSplitFrRulePLow{} to \WidenSplitFrRulePHigh{}), \textsf{M} \WidenSplitFrRuleMMean{}\%; encoder \textsf{P} \WidenSplitFrEncPMean{}\% (\WidenSplitFrEncPLow{} to \WidenSplitFrEncPHigh{}), \textsf{M} \WidenSplitFrEncMMean{}\%.
\emph{Learned.} The advantage on every measure is a property of the original split, not of the setting. The encoder zone is a variant with costs (a second model at several prefix lengths per input, no edit bound, no coverage promise across pivots, no certificate for the free development rewordings), not an improvement we have shown. \cert{} is unchanged.

\FloatBarrier
\section{Do later tasks re-route reworded inputs?}
\label{app:reroute}

\cert{} unchanged, routes only, T5-base word-vectors, the same \WidenTasksEval{} evaluation tasks in the same order of arrival. Inputs: paraphrases of every task (up to three per task and pivot; \RerouteInputs{} for German and French together, \RerouteInputsPivot{} for Spanish) and unedited inputs (up to six per task). For every input and rule the route is recorded right after the input's own task has arrived (only tasks $1,\dots,j$ and their radii used) and after all \WidenTasksEval{} tasks. ``Changed'': the two differ. ``Lost'': own task at arrival, another at the end. The script first reproduced the stored numbers of \cert{} and checked the theorem: no input inside a zone at arrival changed its route under \cert{}.

\paragraph{Lines fixed before each computation.} (1) \emph{First reading.} The zones add something at scale if \nearest{} changes at least 5\% of the routes and \cert{} at most half as many; little if \nearest{} changes fewer than \RerouteLimitLow{}\% or \cert{} more than three quarters as many. Expected: \RerouteExpectNearLow{}\% to \RerouteExpectNearHigh{}\% for \nearest{}, about half of that for \cert{}, all among uncertified inputs, and more inputs to their own task by \nearest{} than by \cert{}. (2) \emph{Follow-up}, added after the first result: the plain nearest signature without radii; ``the radius fixed at arrival is what keeps routes stable'' if it changes at least 5\%, ``rare under any best-match rule'' if fewer than \RerouteLimitLow{}\%, else ``some''. Expected \RerouteExpectSigLow{}\% to \RerouteExpectSigHigh{}\%. (3) \emph{Last rule}, designed after the numbers above on data already looked at, no free setting: zones first (same certificates, same theorem), then the plain nearest signature. Expected 84\% to 88\% to their own task, 1.4\% to 2.0\% changed routes, and on the benchmark no better than \cert{} under heavy rewording. Not an independent test.

\paragraph{Result of the first reading} (Table~\ref{tab:reroute}). All changed inputs lie outside every zone (\nearest{} changed inputs in \RerouteNearTasks{} tasks; the variant \exactfirst{} changed the same). Of the \RerouteCertOwn{} inputs that \cert{} certifies for their own task, \nearest{} gives every one to the same task. Own task at arrival and at the end: \RerouteCertOwnStart{}\% and \RerouteCertOwnEnd{}\% for both rules. Spanish pivot: \RerouteCertEsNChanged{} of \RerouteInputsPivot{} changed under each of the three rules. Unedited inputs: none changed under any rule; \nearest{} serves more of them than \cert{} (\RerouteNearUnOwn{}\% against \RerouteCertUnOwn{}\%). Verdict: ``the zones add little here''. Both expectations about \nearest{} were wrong (the rate, and more inputs to their own task); that all changes are among uncertified inputs was right.

\paragraph{Result of the follow-up and of the last rule} (Table~\ref{tab:reroute}). The plain nearest signature: \RerouteSigNChanged{} changed routes in \RerouteSigTasks{} tasks (verdict ``some''; more was expected), \RerouteSigNLost{} lost inputs, \RerouteSigOwnEnd{}\% to their own task at the end (\RerouteSigOwnStart{}\% at arrival). Of the \RerouteCertOwn{} certified inputs it loses \RerouteSigLostOfCert{} (\RerouteSigLostOfCertPct{}\%). Of its \RerouteSigNChanged{} changed inputs, \RerouteSigToSibling{} end at a task whose definition shares at least 80\% of its words with their own; for \nearest{}, \RerouteNearToSibling{} of \RerouteNearNChanged{}. Spanish pivot: own task \RerouteSigEsOwn{}\% and \RerouteSigEsChanged{}\% changed for the plain rule, \RerouteCertEsOwn{}\% and \RerouteCertEsChanged{}\% for \cert{}, \RerouteNearEsOwn{}\% own for \nearest{}, \RerouteBgEsOwn{}\% own for \scaled{}. \hybrid{} (\RerouteHybNChanged{} changed inputs in \RerouteHybTasks{} tasks, \RerouteHybNLost{} lost) keeps all certificates (\RerouteHybInZone{}\% in a zone) and changes none of the certified inputs; its expectations (84\% to 88\%, 1.4\% to 2.0\%) were right; on the benchmark the plain fallback was no better than \cert{}, as expected, but far worse than thought (\RerouteHybBenchFiveHeavyOwn{}\% under heavy rewording against \RerouteCertBenchFiveHeavyOwn{}\%). Plain-signature re-routes (three of 15, evenly spaced in the file): a \RerouteExOnePivot{} paraphrase of \texttt{\RerouteExOneTask{}} ends at \texttt{\RerouteExOneEnd{}} (\RerouteExOneBetween{} tasks between; \RerouteExOneShare{}\% shared words); a \RerouteExTwoPivot{} paraphrase of \texttt{\RerouteExTwoTask{}} ends at \texttt{\RerouteExTwoEnd{}} (\RerouteExTwoShare{}\%, \RerouteExTwoBetween{}); a \RerouteExThreePivot{} paraphrase of \texttt{\RerouteExThreeTask{}} ends at \texttt{\RerouteExThreeEnd{}} (\RerouteExThreeShare{}\%, \RerouteExThreeBetween{}).

\paragraph{The benchmark tasks} (Table~\ref{tab:reroutebench}). Five tasks, four wordings, routes only, \RerouteBenchInputs{} inputs per wording. \cert{} and \nearest{} give \RerouteCertBenchFiveMediumOwn{}\% to their own task under medium rewording and change \RerouteCertBenchFiveMediumChanged{} and \RerouteCertBenchFiveHeavyChanged{} routes (medium, heavy); the plain signature and \hybrid{} give fewer under both wordings and change many routes (\RerouteSigBenchFiveMediumChanged{} medium; of the \RerouteBenchChangeable{} routes that later tasks could change, \hybrid{} changes \RerouteHybBenchFiveHeavyChanged{} under heavy rewording); on 15 tasks the medium wording gives \RerouteSigBenchFifteenMediumOwn{}\% to the plain signature (\RerouteSigBenchFifteenMediumChanged{} changed). With the original and light wording all four rules give 100\% and change nothing. The plain signature is the better fallback on real tasks with machine paraphrases and much the worse one on the benchmark with heavy development rewordings; we did not test why.

\paragraph{Two fallbacks scaled differently (written down after the results above).} \emph{Reading, not tested.} The radius (a) puts tasks with openings of different lengths on one scale (an average over many tokens lies closer to everything), which the benchmark needs, and (b) shrinks the zone of a task that arrives next to a sibling, which goes wrong at scale. A scale depending only on the length of the opening should keep (a) and drop (b).
\emph{Tried: a fallback scaled by the length of the opening.} With $\sigma(n)$ the root-mean-square distance of the 4{,}000 unrelated inputs from their own mean in the feature $\psi_n$, an input goes to the task with the smallest distance to its signature divided by $\sigma(n_k)$ (ties to the older task); a second version puts the zones of \cert{} in front. No free setting. Test: the benchmark with four wordings (the development set), and paraphrases through Italian, made for this test and the next one and read once, for both verdicts at scale together.
\emph{Line fixed before the run:} success if at least \RerouteLineHeavy{}\% of inputs go to their own task under heavy rewording on five and on fifteen tasks, at least \RerouteLineMedium{}\% under medium, and at least \RerouteLineItalian{}\% on the Italian pivot; against if heavy is below \RerouteLineHeavyAgainst{}\% in either sequence or Italian below \RerouteLineItalianAgainst{}\%; else partial. Expected on the Italian pivot \RerouteExpectNormItLow{}\% to \RerouteExpectNormItHigh{}\% for the length-scaled fallback, \RerouteExpectCfThreeItLow{}\% to \RerouteExpectCfThreeItHigh{}\% with zones in front, and \RerouteExpectNormChangedLow{}\% to \RerouteExpectNormChangedHigh{}\% changed routes.
\emph{Verdict: against.} Heavy: \RerouteNormHeavyFive{}\% (five tasks) and \RerouteNormHeavyFifteen{}\% (fifteen); medium: \RerouteNormMediumFive{}\% and \RerouteNormMediumFifteen{}\%; Italian pivot \RerouteNormItOwn{}\% (\RerouteCertThreeItOwn{}\% with zones in front; \RerouteNormItChanged{}\% changed routes). The Italian pivot was slightly above expectation, the benchmark failed the line, and the changed routes were in the expected range. The reading was only partly right: the correlation of $\sigma(n_k)$ with the radius over the \WidenTasksEval{} tasks is \RerouteSigmaCorr{} (Spearman \RerouteSigmaCorrSpearman{}), so the radius is not mainly a scale for the length of the opening. French and German pivot: \RerouteNormOwnEnd{}\% (\RerouteCertThreeOwnEnd{}\% with zones), \RerouteNormChanged{}\% changed.

\emph{Second reading.} The radius is half the distance to the nearest \emph{known} signature (unrelated or earlier), so it measures how isolated an instruction is. On the benchmark $U_k$ ranges from \RerouteBBenchFiveMin{} to \RerouteBBenchFiveMax{} on the five tasks, and the radius is half of it for \RerouteBBenchFiveShare{}\% of them and for \RerouteBBenchFifteenCount{} of the fifteen. At scale $U_k$ has median \RerouteBMedian{} (range \RerouteBMin{} to \RerouteBMax{}), its correlation with the radius is \RerouteBCorr{}, and the radius equals half of it for only \RerouteBShare{}\% of the \WidenTasksEval{} tasks: for the others an earlier task is nearer than every unrelated one, the radius shrinks, and the nearest-zone fallback turns away exactly their reworded inputs.
\emph{Tried: \scaled{}.} An input goes to the task with the smallest distance divided by $U_k$, fixed when task $k$ arrives (ties to the older task); \certscaled{} puts the zones of \cert{} in front, same certificates. Test and lines as for the first. This rule and its lines were written down while the first test was still running, before any number of the Italian pivot had been read; the Italian numbers of both rules were then read together. Expected \RerouteExpectBgHeavyLow{}\% to \RerouteExpectBgHeavyHigh{}\% under heavy rewording and \RerouteExpectBgItLow{}\% to \RerouteExpectBgItHigh{}\% on the Italian pivot for \scaled{}, \RerouteExpectCfItLow{}\% to \RerouteExpectCfItHigh{}\% for \certscaled{}.
\emph{Verdict: success} (Tables~\ref{tab:reroute}, \ref{tab:reroutebench}). \scaled{} equals the fallback of \cert{} on the benchmark (\RerouteBgHeavyFive{}\% heavy, \RerouteBgMediumFive{}\% medium) and gives \RerouteBgItOwn{}\% on the Italian pivot. \certscaled{} gives \RerouteCertFourItOwn{}\% Italian, \RerouteCertFourOwnEnd{}\% French and German and \RerouteCertFourEsOwn{}\% Spanish, against \RerouteCertItOwn{}\%, \RerouteCertOwnEnd{}\% and \RerouteCertEsOwn{}\% for \cert{}; \RerouteCertFourInZone{}\% of French and German inputs are inside a zone at arrival and none changed, and on the benchmark the shares are those of \cert{} in every wording. Without zones \scaled{} gives \RerouteBgOwnEnd{}\% (French and German); unedited inputs \RerouteBgUnOwn{}\%, against \RerouteCertFourUnOwn{}\% with the zones, because the zones and their collisions are the same as in \cert{}. Less was expected on the 700 tasks (Italian: \RerouteExpectBgItLow{}\% to \RerouteExpectBgItHigh{}\% and \RerouteExpectCfItLow{}\% to \RerouteExpectCfItHigh{}\%); both are higher. The benchmark expectation was right.
\emph{Identity on the benchmark, input by input.} For every benchmark input and wording the route of \cert{} was compared with that of \certscaled{} (and \nearest{} with \scaled{}) after the last task and at every point of the sequence. Five tasks: the radius is half of $U_k$ for all five, so the two fallbacks order the tasks identically; \RerouteIdFiveLast{} differing routes after the last task and \RerouteIdFiveScored{} at any point, in \RerouteIdFiveComparisons{} comparisons (four wordings). Fifteen tasks: the radius is half of $U_k$ for all tasks except \texttt{\RerouteIdTask{}}, the \RerouteIdTaskPos{}rd to arrive, whose radius (\RerouteIdRadius{}) is half the distance to the earlier task \texttt{\RerouteIdEarlier{}} and slightly below half of $U_k$ (\RerouteIdHalfB{}); \RerouteIdFifteenLast{} routes differ after the last task and \RerouteIdFifteenScored{} at any point from the arrival of the input's own task on, in every wording; before their own task has arrived \RerouteIdFifteenBeforeOriginal{} (original), \RerouteIdFifteenBeforeLight{} (light), \RerouteIdFifteenBeforeMedium{} (medium) and \RerouteIdFifteenBeforeHeavy{} (heavy) inputs differ, where the GPU runs take no score. Hence every GPU score of \cert{} and \nearest{} in this paper, for five and fifteen tasks, is also the score of the variant (the comparison covers the four wordings, the five-task order and fifteen-task order 1, \RerouteBenchInputs{} inputs per wording and \RerouteIdFifteenOriginalN{} for the fifteen-task original wording).
\emph{Costs} (Table~\ref{tab:reroute}). \scaled{} and \certscaled{} change more routes after arrival than \cert{}, \nearest{} and fewer than the plain signature; Spanish pivot: \certscaled{} \RerouteCertFourEsChanged{}\% to \scaled{} \RerouteBgEsChanged{}\%; inputs that lose their own task: \RerouteBgNLost{} (\scaled{}) and \RerouteCertFourNLost{} (\certscaled{}). Italian pivot, other rules: \cert{} \RerouteCertItOwn{}\% (\textsf{P} \WidenRefTfiveItP{}\%, \textsf{M} \WidenRefTfiveItM{}\%, \textsf{R} \WidenRefTfiveItR{}\%), zones then plain \RerouteHybItOwn{}\%, length-scaled \RerouteNormItOwn{}\%; changed routes on Italian: \RerouteBgItChanged{}\% (\scaled{}), \RerouteCertFourItChanged{}\% (\certscaled{}), \RerouteCertItChanged{}\% (\cert{}), \RerouteNearItChanged{}\% (\nearest{}), \RerouteSigItChanged{}\% (plain signature).

\begin{table}[t]
  \caption{The stress test of Section~\ref{sec:stress} with the scaled fallbacks: a sixth task whose definition is the light (upper block) or the medium (lower block) rewording of the definition of task~2 arrives after the five tasks; old test inputs of task~2 (\RerouteStressOld{} per wording) routed elsewhere afterwards. Original and light inputs are certified under every rule with zones. In the lower block the medium wording is the look-alike's own wording, so its \RerouteStressOld{} are re-routed under every rule.}
  \label{tab:reroutestress}
  \centering
  \footnotesize
  \begin{tabular}{lrrrr}
    \toprule
    Rule & \shortstack[r]{Routes\\changed,\\original $\downarrow$} & \shortstack[r]{Routes\\changed,\\light $\downarrow$} & \shortstack[r]{Routes\\changed,\\medium $\downarrow$} & \shortstack[r]{Routes\\changed,\\heavy $\downarrow$} \\
    \midrule
    \multicolumn{5}{l}{New definition: the light rewording} \\
    \cert{} & 0 & 0 & 0 & 0 \\
    \nearest{} & 0 & 100 & 0 & 0 \\
    Length-scaled fallback & 0 & 100 & 78 & 23 \\
    Zones, then length-scaled & 0 & 0 & 78 & 23 \\
    \scaled{} & 0 & 100 & 2 & 20 \\
    \certscaled{} & 0 & 0 & 2 & 20 \\
    \midrule
    \multicolumn{5}{l}{New definition: the medium rewording} \\
    \cert{} & 0 & 0 & 100 & 2 \\
    \nearest{} & 0 & 0 & 100 & 2 \\
    Length-scaled fallback & 0 & 0 & 100 & 1 \\
    Zones, then length-scaled & 0 & 0 & 100 & 1 \\
    \scaled{} & 0 & 0 & 100 & 80 \\
    \certscaled{} & 0 & 0 & 100 & 80 \\
    \bottomrule
\end{tabular}

\end{table}

Stress test: with the light rewording as the look-alike's definition, the look-alike takes \RerouteStressTaken{} uncertified old inputs under \certscaled{} (\RerouteStressCertFourMedium{} medium, \RerouteStressCertFourHeavy{} heavy) that the fallback of \cert{} kept (0 and 0); with the medium rewording as its definition it takes \RerouteStressMedCertFourHeavy{} of the \RerouteStressOld{} old inputs with the heavy wording, where \cert{} loses \RerouteStressMedCertHeavy{}, so \RerouteStressMedExtra{} more; the \RerouteStressOld{} certified old inputs of the original and light wording do not change under any rule with zones. The length-scaled fallback changes \RerouteStressNormMedium{} and \RerouteStressNormHeavy{}.

\paragraph{What the tests show and do not show.} The fallback \scaled{} was found at the end, with its rule fixed in advance and its verdict read on paraphrases not read before; zones, certificates and theorem are untouched. It is a variant: every GPU run of \cert{} and every other route test used the fallback of \cert{}; the benchmark rewordings are the development set, so there the test is not independent; the reasons for the success are read from the numbers, not tested. With zones in front the best rule without zones still gives \RerouteZonesCost{} (Italian), \RerouteZonesCostFG{} (French and German) and \RerouteZonesCostEs{} (Spanish) points more right routes: the zones add the proof, not better routing. At this scale the certified switch showed no measured stability beyond the nearest-zone rule: both changed the same inputs. What it adds is the proof for the inputs inside a zone, about half of the machine-paraphrased inputs. The first reading went against the claim that the zones prevent re-routing a best-match rule would cause; in this collection that case is rare. The plain signature serves more reworded inputs than \cert{} and pays with about two silent changes in a hundred; \hybrid{} gains \RerouteHybGain{} points in one test and loses \RerouteHybLoss{} in the other. The paraphrases are made by translation models and the collection is one collection.
